\documentclass[11pt,a4paper]{article}

\usepackage[T1]{fontenc}
\usepackage{libertinus}
\usepackage[a4paper,margin=2.4cm]{geometry}
\usepackage{microtype}
\usepackage{booktabs,tabularx,array}
\usepackage{graphicx}
\usepackage[labelfont=bf,font=small,skip=4pt]{caption}
\usepackage{subcaption}
\usepackage{float}
\usepackage{enumitem}
\usepackage{xcolor}
\usepackage[hidelinks]{hyperref}

\graphicspath{{figures/}}
\setlist{itemsep=2pt,topsep=4pt}
\setlength{\parindent}{0pt}
\setlength{\parskip}{5pt}

% Class colours: the same as the boxes in the figures (render.py)
\definecolor{c-person}{RGB}{31,119,180}
\definecolor{c-cyclist}{RGB}{44,160,44}
\definecolor{c-car}{RGB}{214,39,40}
\definecolor{c-escooter}{RGB}{148,103,189}
\definecolor{c-suv}{RGB}{255,127,14}
\definecolor{c-motorcyclist}{RGB}{140,86,75}
\definecolor{c-bus}{RGB}{227,119,194}
\definecolor{c-van}{RGB}{188,189,34}
\definecolor{c-truck}{RGB}{127,127,127}
\newcommand{\sw}[1]{\textcolor{#1}{\rule[-0.1ex]{0.75em}{0.75em}}}
\newcommand{\cls}[1]{\texttt{#1}}

% One image panel of a figure: \panel{width}{file}{caption}
\newcommand{\panel}[3]{%
  \begin{subfigure}[t]{#1\linewidth}\centering
    \includegraphics[width=\linewidth,height=4.3cm,keepaspectratio]{#2}
    \caption{#3}
  \end{subfigure}}

\begin{document}

{\LARGE\bfseries CAMINA labelling guide}\par\vspace{2pt}
{\large Nine road-user classes for a roadside counting sensor}\par\vspace{4pt}
{\small\color{black!60} UCD Spatial Dynamics Lab \quad·\quad Version 1.0, 25 September 2026}
\vspace{10pt}

CAMINA counts the road users that cross a line in a camera's view. The detector learns
exactly what the labels say, so one consistent rule matters more than many images. This
guide fixes those rules: for correcting the current dataset and for labelling new images
from Dublin cameras.

The current dataset (TRA 2026: 1,837 images, 13,148 boxes) was labelled by an automatic
pipeline. A COCO detector supplied the common classes, an open-vocabulary detector
(YOLO-World) supplied \cls{e-scooter}, \cls{SUV} and \cls{delivery\_van}, and cyclist and
e-scooter boxes were built by joining a person to the vehicle beneath them. How many labels
were then corrected by hand is not recorded. Every example below is from that dataset, shown
with its labels as they are; Figure~\ref{fig:errors} shows errors that remain.

\section*{Five rules}

\begin{enumerate}[leftmargin=*,label=\textbf{\arabic*.}]
  \item \textbf{Label every road user in the image.} \emph{Why:} an unlabelled road user
    teaches the model that it is background.
  \item \textbf{One box per road user.} A rider and the bicycle, e-scooter or motorcycle
    they ride share one box; the rider gets no separate \cls{person} box.
    \emph{Why:} the sensor counts one track per road user, and a second box would count the
    rider twice. The dataset already does this for 97\% of cyclists.
  \item \textbf{Box what is visible, tightly.} Do not extend a box behind an occluding
    object. Label objects cut by the image edge. \emph{Why:} the detector can only learn
    what it can see, and guessed extents differ from one annotator to the next.
  \item \textbf{A rider class needs a rider.} A parked bicycle, e-scooter or motorcycle is
    not labelled; a person pushing one is a \cls{person}. A rider stopped with a foot down
    is still a rider. \emph{Why:} the classes describe how someone is travelling; a person
    walking a bicycle moves, and looks, like a pedestrian.
  \item \textbf{When unsure, choose the broader class, or skip.} Between \cls{car} and
    \cls{SUV}, choose \cls{car}. If you cannot tell what an object is, leave it
    unlabelled. \emph{Why:} a wrong label on a rare class costs more than a missed one, and
    an object no one can identify cannot be labelled consistently.
\end{enumerate}

\clearpage
\section*{The nine classes}

{\small
\renewcommand{\arraystretch}{1.25}
\begin{tabularx}{\linewidth}{@{}l X X r@{}}
\toprule
Class & What it is & Not this & Boxes \\
\midrule
\sw{c-person}\ \cls{person} & Anyone on foot, standing or sitting; includes a person
  pushing a bicycle or scooter, on a skateboard, or in a wheelchair
  & People inside vehicles; people in posters or reflections & 6,975 \\
\sw{c-cyclist}\ \cls{cyclist} & A person riding any bicycle, e-bike or cargo bike; one box
  around rider and bicycle & A parked bicycle; a person walking a bicycle & 2,012 \\
\sw{c-car}\ \cls{car} & A passenger car: saloon, hatchback, estate, coupé, MPV, taxi
  & SUVs and crossovers & 2,105 \\
\sw{c-escooter}\ \cls{e-scooter} & A person standing on a kick scooter, electric or not;
  one box around rider and scooter & A parked or carried scooter; a seated moped rider & 728 \\
\sw{c-suv}\ \cls{SUV} & A sport utility vehicle or crossover: a tall body sitting high over
  the wheels, upright tailgate (Toyota RAV4, Nissan Qashqai, Hyundai Tucson)
  & MPVs (\cls{car}); pickups (\cls{truck}) & 456 \\
\sw{c-motorcyclist}\ \cls{motorcyclist} & A person riding a motorcycle or moped; one box
  around rider, passenger and vehicle & A parked motorcycle & 307 \\
\sw{c-bus}\ \cls{bus} & A bus or coach, single- or double-deck & A tram (Luas): not
  labelled & 321 \\
\sw{c-van}\ \cls{delivery\_van} & A van body, whatever it carries (Ford Transit,
  VW Transporter, Renault Kangoo) & A box body behind a separate cab (\cls{truck}) & 112 \\
\sw{c-truck}\ \cls{truck} & A goods vehicle with a separate cab: box truck, lorry, tipper,
  bin lorry, pickup; an articulated lorry is one box & Vans & 132 \\
\bottomrule
\end{tabularx}}

\paragraph{Three choices, and why.}
\emph{Pickups are trucks}, as in COCO, so pre-labels from COCO-trained models need no
correction there. \emph{Trams are not labelled}: there is no tram class, and labelling the
Luas as \cls{bus} would add trams to the bus count. \emph{Any van body is a
\cls{delivery\_van}}: whether a van is delivering cannot be seen, and its shape can.

\begin{figure}[H]
  \centering
  \panel{0.27}{rider_cyclist}{\cls{cyclist}: rider and bicycle in one box.}\hfill
  \panel{0.19}{rider_escooter}{\cls{e-scooter}: rider and scooter.}\hfill
  \panel{0.2}{rider_stopped}{Stopped, one foot down: still \cls{e-scooter}.}\hfill
  \panel{0.27}{rider_motorcyclist}{\cls{motorcyclist}: rider and moped.}
  \caption{Rider classes: one box around the rider and the vehicle (rules 2 and 4).}
  \label{fig:riders}
\end{figure}

\clearpage
\begin{figure}[H]
  \centering
  \panel{0.3}{car_sedan}{\cls{car}: a low saloon body.}\hfill
  \panel{0.2}{suv_4x4}{\cls{SUV}: a tall body, high over the wheels.}\hfill
  \panel{0.44}{suv_street}{Two \cls{SUV}s; the officer is a \cls{person}.}
  \caption{Car or SUV. Look at the height of the body over the wheels; if still unsure,
    choose \cls{car} (rule 5).}
  \label{fig:suv}
\end{figure}

\begin{figure}[H]
  \centering
  \panel{0.24}{van}{\cls{delivery\_van}: a van body.}\hfill
  \panel{0.3}{truck}{\cls{truck}: a box behind a separate cab.}\hfill
  \panel{0.36}{bus}{\cls{bus}.}
  \caption{Vans, trucks and buses.}
  \label{fig:large}
\end{figure}

\clearpage
\begin{figure}[H]
  \centering
  \panel{0.32}{err_parked_bicycles}{Parked bicycles, and a person pushing one, labelled
    \cls{cyclist}. Should be: nothing; \cls{person}.}\hfill
  \panel{0.2}{err_rider_as_person}{An e-scooter rider labelled \cls{person}. Should be:
    \cls{e-scooter}.}\hfill
  \panel{0.4}{err_car_as_truck}{A car on its roof labelled \cls{truck}. Should be:
    \cls{car}.}

  \medskip
  \panel{0.3}{err_tram_as_bus}{A tram labelled \cls{bus}. Should be: nothing.}\qquad
  \panel{0.32}{err_two_classes}{One SUV with two boxes, \cls{SUV} and
    \cls{delivery\_van}. Should be: one \cls{SUV} box.}
  \caption{Errors that remain in the current dataset. Fix these when you meet them.}
  \label{fig:errors}
\end{figure}

\section*{Workflow for new images}

\begin{enumerate}[leftmargin=*,label=\textbf{\arabic*.}]
  \item \textbf{Pre-label} with the most accurate model available; it runs on a GPU, so it
    need not be the Raspberry Pi model.
  \item \textbf{Review every box}: correct its class and its extent, and delete boxes on
    things that are not road users.
  \item \textbf{Look for what is missing}, especially small, distant and partly hidden road
    users. Pre-labels make it easy to accept what is there and forget what is not.
  \item \textbf{Check one class at a time} over the whole batch (all SUVs, then all vans):
    inconsistencies stand out side by side.
  \item \textbf{Keep the camera in the file name}, so the dataset split can keep each
    camera's images together and the test score stays honest.
\end{enumerate}

\vfill
{\small\color{black!60} The figures are generated from the committed labels by
\texttt{docs/labelling\_guide/render.py}; box counts are from
\texttt{training/dataset/split.json}. The class list and its order are fixed by
\texttt{configs/classes.yaml}.}

\end{document}
